\chapter{Hematoma}

\section*{Definition}
A hematoma is a collection of blood outside blood vessels, usually in liquid form within the tissues. It results from trauma that causes blood vessel rupture. Hematomas can occur anywhere in the body and range from minor bruises to life-threatening intracranial or retroperitoneal bleeds.

\section*{What Happens in Your Body}
When a blood vessel is damaged, blood escapes into the surrounding tissue. The blood clot forms a localized collection. Over time, the body breaks down the blood products, leading to color changes (red to purple to green to yellow). Large hematomas can cause compression of adjacent structures, organ dysfunction, or hypovolemia.

\section*{Causes}
\begin{itemize}
\item Trauma (most common): contusions, fractures, crush injuries
\item Surgery (postoperative hematoma)
\item Anticoagulant therapy (warfarin, heparin, DOACs)
\item Coagulation disorders (hemophilia, thrombocytopenia, liver disease)
\item Spontaneous: ruptured aneurysm, vascular malformation
\item Repeated minor trauma (subungual hematoma)
\item Venipuncture or injection
\end{itemize}

\section*{When to See a Doctor}
\begin{itemize}
\item Localized swelling and discoloration
\item Pain and tenderness at the site
\item Firm or fluctuant mass on palpation
\item Ecchymosis (black-and-blue discoloration)
\item Neurological deficits if intracranial
\item Signs of hypovolemia if large or internal
\end{itemize}

\section*{Related Symptoms}
\begin{itemize}
\item Bruise (ecchymosis, superficial hematoma)
\item Subdural hematoma (intracranial)
\item Epidural hematoma (intracranial)
\item Subungual hematoma (under the nail)
\item Retroperitoneal hematoma
\item Hypovolemic shock
\end{itemize}

\section*{Self-Care / Home Management}
\begin{itemize}
\item RICE protocol: Rest, Ice, Compression, Elevation
\item Apply ice for the first 24-48 hours, then warm compresses
\item Elevate the affected area
\item Take acetaminophen for pain (avoid NSAIDs if bleeding risk)
\item Monitor for signs of expansion or infection
\end{itemize}

\section*{How Your Doctor Will Evaluate You}
\begin{itemize}
\item Clinical examination
\item Ultrasound for superficial or deep hematomas
\item CT scan for intracranial, thoracic, or abdominal hematomas
\item MRI for better characterization
\item Complete blood count and coagulation profile
\end{itemize}

\section*{Treatment Options}
\begin{itemize}
\item Observation for small, stable hematomas
\item Drainage for large, tense, or infected hematomas
\item Surgical evacuation for intracranial hematomas causing mass effect
\item Reverse anticoagulation if applicable
\item Compression bandaging for extremity hematomas
\end{itemize}
\section*{References}
\begin{thebibliography}{99}
\bibitem{hematoma:ref1} Trafton PG. ``Hematoma.'' In: Marx JA, Hockberger RS, Walls RM, eds. Rosen's Emergency Medicine. 9th ed. Elsevier; 2018.
\bibitem{hematoma:ref2} Millington JT, Millington S. ``Subdural Hematoma.'' In: Brain Injury Medicine. 2nd ed. Demos Medical; 2012.
\bibitem{hematoma:ref3} Pierce CA, et al. ``Retroperitoneal Hematoma.'' Am J Surg. 2005;189(5):540--543.
\bibitem{hematoma:ref4} Powers WJ, Rabinstein AA, Ackerson T, et al. ``Guidelines for the Early Management of Patients with Acute Ischemic Stroke.'' Stroke. 2019;50(12):e344--e418.
\bibitem{hematoma:ref5} Bullock MR, Chesnut R, Ghajar J, et al. ``Surgical Management of Traumatic Brain Injury.'' Neurosurgery. 2006;58(3 Suppl):S2--1--S2--60.
\end{thebibliography}
